Hello leudde!

Kommendes Wochenende kommt einiges auf uns zu, dazu bitte ich euch schon mal die (vorläufige!) Schichtplanung anzuschauen.

Der letzte Sonntag war wieder ein Beispiel dafür, wie es ist wenn das Cafe komplett ausgelastet ist. 
Aber: Für uns fühlt sich das viel chaotischer an als es bei 90prozent der Gäste ankommt.
So einen Ansturm und Chaos bewältigen zu können ist für Niemanden eine leichte Aufgabe, vorallem wenn Gäste mit dem Gastgarten machen was sie wollen.

Um es aber für uns leichter zu machen und euch zu informieren, möchte ich (für ALLE) noch ein paar wichtige REMARKS/REMINDER für die Zukunft zusammenfassen: 
(Tut mir leid es wird viel)

- zweites Coha Handy wird wieder aktiviert
    (Tap2pay auf meinem Android funktioniert nicht gut, da oft offline geht und Bestellungen verloren gehen)

- neue Wasser Gläser sind bestellt
    
- Bestellungen von Gästen, die noch KEINEN TISCH haben, werden NICHT an der Theke aufgenommen.
    Da dürft und sollt ihr gerne streng sein. Trotzdem nicht sofort wegschicken, sondern höflich darauf verweisen, am Tisch zu bestellen.
    An Theke nur Togo Bestellungen, Kassieren, Fragen beantworten etc.

- Dafür Service 1 bitte SO VIEL WIE MÖGLICH DRAUSSEN sein. Nicht warten bis Gäste rein kommen und dann sagen dass wir eh raus kommen.
    Aber dieses Gefühl von 'Es kommt eh gleich wer' haben sie nicht wenn Service drinnen ist.
    In der Zeit von Hinsetzten bis Bestellen sind sie am ungeduldigsten.
    Service macht Pausen um kurz durchzuschnaufen, was zu trinken, Tische umzubuchen etc. und erst DANN Geschirr, Bar helfen etc.
    Wenn weniger los ist, kann man natürlich ein bisschen flexibler sein

- Dafür wiederrum muss Bar/Küche, dem Service den 'Rücken frei halten'. In Stehzeiten (nach ebenfalls kurz durchschnaufen und was trinken)
    --> Bestände checken (Geht bald was aus?), Service informieren über zB letzte Bagels/letztes Stk Kuchen, dreckiges Geschirr so schnell wie möglich waschen, etc.
    Backend ist genauso wichtig wie schöne cappuccini machen.
    Service ist einer der herausfordernsten Aufgaben und muss sich daher auf Bar und Backend verlassen können.

- So wie am Sonntag, is für Service 1 unmöglich alleine alle Gäste mit aufnehmen, Kassieren, Tische säubern zu betreuen. Es werden immer Gäste reinkommen.
    Dann kann an Theke bezahlt werden, aber Bar 1 unterbricht seine Arbeit NIE wenn ein Tisch/Tablett nicht komplett fertig ist, damit Service während Gäste an Theke bedient werden zumindest fertiges Tablett bringen kann. 
    Wenn wirklich laufend Bons kommen, betreut vorzugsweise Bar2 Gäste an Theke damit Getränke Produktion so wenig wie möglich unterbrochen wird.

- Service bitte aktiv Tische abräumen und leere Gläser laufend einsammeln. Bei dieser Gelegenheit immer gerne gleich fragen ob sie noch was trinken wollen.
    Damit verhindert man wieder dass sie von selbst reinkommen und drinnen bestellen wollen

- Für die Tatsache dass es manchmal unter der Woche zu zweit auch fast alles voll und extrem stressig ist, ist uns noch keine Lösung eingefallen, da das sehr unvorhersehbar passiert.
    Als einzigen Tipp kann ich weitergeben: Nicht stressen lassen. Wir haben bis jetzt und werden immer tun, so viel wir können. Einfach höflich auf mögliche Wartezeiten hinweisen und um Verständnis bitten. Es is ganz normal dass wenn 30 leute innerhalb von 15min kommen, nicht alle sofort bedient werden können. Deshalb braucht ihr euch gar nicht schlecht fühlen dafür.
    
Wir sind aber auch jederzeit offen für Feedback/Fragen und möchten dass jede Art von Unklarheiten oder Hinweise eurerseits immer offen mit uns kommuniziert werden können. 

Danke :)